\chapter{Streaming, language, code and persistent agents}
\label{evod49:chapter}
A code agent submits a tool action, the process crashes, and no acknowledgement
was recorded. Should restart submit the action again? The answer depends on
what the external service did, not on how confident the model sounds. This
small ambiguity connects streaming, memory, authority and persistence more
directly than a large prompt does.

We retain the supplied event/session example but replace its in-memory tool
cache with replayable intents and results. The result is a serialized recovery
oracle: it emits instructions to dispatch or reconcile, but executes no tools,
stores no durable log and implements no production concurrency. Language and code
are application workloads requiring their own trained models and evaluators.
Neither the ACGT Transformer nor the restricted native recurrence becomes a
language agent by being placed behind this shell.

\section{Events, ordering and bounded pressure}
\index{backpressure}\index{event ordering}\index{immutable payloads}
\index{persistent agents!event log}
An event is $(i,k,p)$: a sequence number $i$, a kind $k$ and a typed payload
$p$. Within one session the log orders accepted observations and decisions.
It does not recover the true wall-clock order of all outside events. If two
producers race, a sequencer or transactional append rule must decide the order.
The reference has one serialized caller, so $i$ is exactly the previous length.

User input, proposals and emitted output are observations. Tool intents and
results change a recovery state machine. Closure is a terminal event. Large
artifacts belong behind immutable handles with content identity and access
checks; copying binaries into every event multiplies retention and privacy cost.
Our small payloads are canonical JSON strings. Freezing only a dataclass while
retaining a mutable dictionary would still permit old history to change through
an alias.

Let $a_t$ be arrivals, $d_t$ departures and $q_t$ queue occupancy after service
at tick $t$. For a capacity $K$, accepted arrivals satisfy
\[
q_{t+1}=q_t+a_t-d_t,\qquad 0\le q_{t+1}\le K.
\]
If an arrival would violate $K$, the protocol must block, reject, persist to
another bounded tier or shed under a declared policy. Changing $K$ cannot solve
a sustained arrival rate above service rate. Dropping a tool completion is also
different from dropping a best-effort telemetry sample.
\Listing{Stream}
The generated FIFO test pushes two items, rejects the third, then returns
the first two unchanged. Its conservation check concerns accepted items only;
the rejected producer remains responsible for the item.
Capacity here counts items, not bytes; large variable payloads require a byte
budget too. The reference retains its event list in memory without compaction.
It establishes neither bounded total-history storage nor durable retention.

\Plate{01-pressure}{The upper trace shows bounded queue occupancy; the third
arrival is rejected while two credits are occupied. The lower trace shows
the two accepted items leaving in FIFO order. Rejection transfers no ownership
and must be visible to the producer.}{evod49:pressure}

\section{Four states that cannot substitute for one another}
\index{memory!tiers}\index{native state}\index{KV cache}
Keep engine-native state, active model context, durable addressable memory
and authoritative external state separate. DOGMA-style native state is a
model-specific compressed representation. Hermon's per-layer K/V retains
addressable attention history. Active context is the input selected for a
model call. Durable memory is a versioned record store. Files, databases and
tool-side objects are authoritative external state.

MemGPT provides prior work on managing a bounded context through memory tiers
and function calls \cite{evod49:memgpt}. The applicable systems lesson is
explicit movement between tiers. A summary remains a lossy view of source
events. It cannot silently become an instruction with greater authority than
its source, and a deletion/retention policy must address summaries and indexes
as well as original records.

The event log is not automatically an engine checkpoint. To resume Chapter 31,
the native carry needs its artifact, owner and consumed-token position. For
Chapter 35, the ordered K/V at every layer needs its model identity and learned
absolute positions. Their shapes and meanings differ. The local tests call
the real stable interfaces and check whole-versus-chunk continuation on their
actual vocabularies. No conversion between the two state types is invented.

\section{Identity binds replay to one contract}
\index{session identity}\index{checkpoint!prefix digest}
A session identity contains tenant, session, artifact and policy version.
The checkpoint also binds context limit, action budget and lease expiration.
These latter fields are contract configuration, not text that the model may
rewrite. Migration to a new model or policy requires an explicit transition
and an appropriate state reset or proven migration.

For a log prefix of $n$ events define
\begin{equation}\label{evod49:digest}
D_n=\operatorname{SHA256}(\operatorname{canonical}(I,C,E_{0:n})),
\end{equation}
where $I$ is identity and $C$ the session contract. We use an event count $n$,
so an empty prefix has $n=0$. This avoids interpreting a negative array index
as a checkpoint. Later append-only events leave $D_n$ unchanged. Altering a
prefix payload, artifact, policy or budget invalidates the checkpoint.

A digest detects inconsistency against a trusted recorded digest. An adversary
able to replace both log and checkpoint can recompute both. Authentication,
authorization, encryption, retention and crash-atomic storage remain separate
requirements. A valid digest also does not prove that a recorded tool result
is true; it only binds the bytes of that observation.

\section{An action state machine for ambiguous crashes}
\index{idempotency key}\index{reconciliation}\index{budget reservation}
\index{idempotency!recovery}
For each operation key $k$, record a fingerprinted command and one of three
states: absent, intent-only, or completed. The command includes operation name
and canonical arguments. Reusing $k$ with a different command is an error.
The key's scope includes session/tenant identity; global raw keys would allow
cross-session conflicts.

\Plate{02-recovery}{The same operation key can denote only one command. A new
admission records intent before dispatch. A crash before a result leaves an
unresolved intent, which requires reconciliation. Only a recorded completion
allows returning a cached result. No transition assumes that a timeout means
the outside effect did not occur.}{evod49:recovery}

\Listing{Session.issue}
The policy decision is supplied by trusted application code, never extracted
from a proposal string. The reference requires the literal Boolean true,
rejects an expired lease, then reserves one action budget unit before dispatch.
Repeated lookups do not reserve another unit. Admission counts distinct intents,
including failed or unresolved actions; refund policies would be a separate
contract. An action admitted before expiration may complete afterward, so its
observation must still be recorded. Expiration blocks new dispatch, not truth.
Ticks and authorization decisions are trusted caller inputs in this oracle;
a host must supply a trusted clock and policy service. The code does not read
an external clock or protect itself against a privileged caller forging history.

\Listing{Session.complete}
Duplicate identical completion is harmless; a different value for the same key
is a conflict. Closure requires every intent to have a result in this narrow
reference. A real service might instead close with explicit unresolved/cancelled
states and a separate reconciliation owner. It cannot silently forget them.

\begin{proposition}[At most one local dispatch instruction per key]
In one serialized valid session history, the reference returns
\texttt{dispatch} for a key at most once.
\end{proposition}
\begin{proof}
The first dispatch records an intent for that key before returning. Every
later issue either rejects changed arguments, returns the recorded result,
or returns reconciliation. None of those paths dispatches. Resume reconstructs
the same intent set from the log before allowing further issue calls.
\end{proof}
The theorem is about local instructions. It is not exactly-once external
execution. A dispatcher can fail between accepting the instruction and sending
it, or an external tool can execute before the result is logged. A tool protocol
with durable idempotency may permit a same-key retry; without it, query the
authoritative state or require reconciliation by an operator. No prompt can
supply missing external transactional semantics.

\begin{center}\input{\Generated/trace.tex}\end{center}
This generated trace has one admitted read, restart with no result, then a
recorded result. It is a constructed event sequence, not a failure-injected
distributed-service measurement.

\section{Replay is reduction, not re-execution}
\index{replay}\index{crash recovery}
\Listing{resume}
Replay applies validated event transitions to fresh derived state and verifies
the checkpoint prefix. The full provided suffix is also validated. It does
not call a model, rerun a shell command or resend an action. Its event schemas
reject out-of-order sequence numbers, result-before-intent, duplicate intents,
and events after closure. The log is trusted input for authorization history:
replay validation establishes internal consistency, not that its original
producer had permission.

Use induction to justify reconstruction. The empty log yields empty intent
and result maps. A valid observation leaves them unchanged. A valid intent
adds exactly the recorded key/command; a valid result adds its unique value.
Thus after every prefix the replay maps equal the maps obtained by original
serialized application of that prefix. This proves replay consistency under
deterministic reduction, not disk durability or linearizability.

Model replay is a different issue. Reproducing sampled tokens additionally
requires the artifact, sampler, RNG state, tokenizer and numerical environment.
For audit, recording emitted tokens and tool observations is usually stronger
than assuming a future call will regenerate them exactly. Once output has been
emitted, a rollback cannot make a recipient unsee it.

\section{Language and code workloads need task evidence}
\index{agent-computer interface}\index{sandbox}\index{patch validation}
For language streams, measure turn completion, contextual consistency,
retrieval quality, latency and cancellation under bounded context. Retain
authoritative user constraints outside an evictable summary. Test whether an
old correction survives summarization, whether a deleted source is still
retrieved, and whether ambiguous references produce calibrated deferral.
These are evaluation proposals, not capabilities demonstrated by our event log.

For a coding task bind repository revision, working tree identity, task,
dependency lockfiles, environment and allowed operations. A patch candidate
is meaningful only relative to that base. Inspection, editing, static checks,
targeted tests and diff review provide different evidence. A passed unit test
does not certify security, compatibility or every hidden requirement.

SWE-agent studies how an agent-computer interface changes software-engineering
agent performance \cite{evod49:swe}. Our proposed interface uses typed
operation names, structured arguments, immutable resource handles and explicit
errors. It exposes useful observations but never grants authority merely
because a tool name exists. Sandbox filesystem, network, process, credential
and resource permissions are supplied by the host. Generated code and retrieved
repository text are untrusted inputs.

\Plate{03-code}{A patch is evaluated against a specific base revision. If the
base changes during evaluation, the tested candidate and promotion target no
longer agree. The conflict path requires rebase/revalidation rather than reusing
the earlier test result. This is a version contract, not an implementation of
repository transactions or a sandbox.}{evod49:code}

\section{Plans, leases, humans and observability}
\index{lease}\index{human assistance}\index{observability}
A long-horizon plan is a versioned hypothesis with preconditions and completion
evidence. Re-read mutable external state before acting on stale assumptions.
A lease specifies lifetime; budgets specify tokens, compute, storage and tools
in their own units. An integer tool-count bound in our code is not monetary
accounting. For concurrent work, reserve resources atomically before launch and
reconcile completion charges rather than checking only after consumption.

If approval is required by the application policy, bind it to the concrete
action, resource version, scope and expiry. A changed action invalidates the
approval. A summary saying ``approved'' is not an approval record. Human
intervention rate, success rate and violations should be reported separately;
silently accepting more risk can increase apparent autonomy.

Trace session, plan, engine and tool spans with durations, queue pressure,
retries, unresolved actions and budget use. Avoid secrets and unnecessary
private content. Concurrent branches may share immutable sources but need a
conflict rule for writes. The absent Chapter 39/40/45 production adapters remain
integration dependencies: this chapter does not assume those contracts compose automatically or
claim a passing multi-tenant persistent-agent deployment.

\section{Selective glossary}
\begin{description}
\item[Intent.] A recorded admission for one fingerprinted operation key, before
its dispatch outcome is known.
\item[Reconciliation.] Determining an unresolved external outcome using its
tool protocol or authoritative state rather than blindly repeating the effect.
\item[Idempotency key.] A scoped identity whose repeated use denotes the same
command; safety depends on the external service honoring the protocol.
\item[Replay.] Deterministic reduction of recorded events into derived state,
separate from re-execution of effects.
\item[Lease.] A bounded lifetime for authority, distinct from a count or cost
budget; late observations may still require recording.
\end{description}

\section{Exercises and worked solutions}
\begin{enumerate}
\item A producer sends three items into an empty capacity-two queue without
service. \textbf{Answer:} two are accepted; the third is rejected and retained
by the producer under this contract.
\item Derive queue conservation. \textbf{Answer:} add accepted arrivals and
subtract departures; neither blocked arrivals nor retries count as new stored items.
\item Can a larger queue fix indefinite overload? \textbf{Answer:} no; it only
delays exhaustion when arrival rate exceeds service rate.
\item Is K/V durable memory? \textbf{Answer:} it is artifact-specific native
attention history, not a general authoritative record store.
\item What does an empty checkpoint count? \textbf{Answer:} zero events,
with a digest over empty history and the bound identity/contract.
\item Append after a checkpoint: must it fail? \textbf{Answer:} no; the bound
prefix remains unchanged and suffix transitions are validated.
\item Can replacing log and digest evade detection? \textbf{Answer:} yes;
unauthenticated hashes do not establish origin or authority.
\item Why canonical strings rather than frozen dict fields? \textbf{Answer:}
nested mutable dictionaries would preserve aliases into history.
\item Same key, changed path? \textbf{Answer:} reject the fingerprint conflict.
\item Crash after intent, before acknowledgement? \textbf{Answer:} reconcile
the external outcome; do not assume absent acknowledgement means absent effect.
\item Does the dispatch theorem prove exactly-once effects? \textbf{Answer:}
no; dispatch and external execution are different boundaries.
\item Does a cached lookup bypass current policy? \textbf{Answer:} not here;
issue checks current authorization and lease before returning a result.
\item Why reserve before dispatch? \textbf{Answer:} post hoc accounting can
overspend when multiple actions are in flight.
\item May an admitted result arrive after expiry? \textbf{Answer:} yes;
recording its outcome is needed for reconciliation even when new dispatch is denied.
\item Replay calls the tool again? \textbf{Answer:} no; it reduces observations
into derived maps and leaves unresolved intents explicit.
\item What does a source-less summary lose? \textbf{Answer:} provenance,
fine detail and potentially valid deletion/authority interpretation.
\item A patch passed tests on revision A; deploy to changed B? \textbf{Answer:}
reconcile the base and revalidate; the old evidence concerns a different object.
\item Design a crash test. \textbf{Rubric:} enumerate before-intent,
after-intent, after-effect and after-result crashes; specify durable records,
tool idempotency, reconciliation and duplicate-result expectations.
\item Design a language-memory test. \textbf{Rubric:} held-out sessions,
provenance/deletion controls, fixed retrieval budget, paired summarization baseline,
success/latency/deferral metrics and audited human intervention.
\item Why are the real-engine continuation tests limited? \textbf{Answer:}
they preserve actual native semantics on existing vocabularies; they do not
train language models, implement cross-engine conversion or validate a service.
\end{enumerate}

\section{Bridge}
Persistent state changes what evidence is available at the next decision.
Chapter 50 asks which changes are learning, how they interfere with old skills,
and what tests an update must face before becoming an active artifact.
